\documentclass[10pt,a4paper]{amsart}
\usepackage[margin=19mm]{geometry}
\usepackage{amsmath,amssymb,mathtools}
\usepackage[scr=rsfs,cal=euler]{mathalfa}
\usepackage{xfrac}
\usepackage[unicode,hidelinks,bookmarks=false]{hyperref}
\newcommand{\ie}{i.e.\ }
\newcommand{\cf}{cf.\ }
\newcommand{\resp}{respectively\ }
\newcommand{\Subsection}{\subsection}
\providecommand{\nobreakdash}{\nobreak\hbox{-}}
\setlength{\emergencystretch}{2em}
\multlinegap0pt
\usepackage{bm}
\usepackage{bbm}
\usepackage{dsfont}
\usepackage{xspace}
\usepackage{cancel}
\usepackage{mathtools}
\usepackage{amsthm}
\makeatletter
\@ifundefined{theorem}{\newtheorem{theorem}{Theorem}}{}
\@ifundefined{corollary}{\newtheorem{corollary}[theorem]{Corollary}}{}
\@ifundefined{lemma}{\newtheorem{lemma}[theorem]{Lemma}}{}
\@ifundefined{proposition}{\newtheorem{proposition}[theorem]{Proposition}}{}
\makeatother
\title[Parseval-type identities for Gowers uniformity norms in fini]{Parseval-type identities for Gowers uniformity norms in finite abelian groups}
\author{Martin Niepel}
\date{Source: arXiv:2508.08968v1 (CC BY 4.0)}
\begin{document}
\maketitle

\noindent\emph{Source and licence.} Parseval-type identities for Gowers uniformity norms in finite abelian groups. Version arXiv:2508.08968v1 (CC BY 4.0), distributed under the Creative Commons Attribution 4.0 International licence, as recorded in the arXiv licence metadata for this exact version and shown on the abstract page https://arxiv.org/abs/2508.08968v1. Full rights notice: licenses/arxiv-2508.08968-CC-BY-4.0.txt.\par\smallskip

\noindent\textit{Editorial scope.} Counted unit: the Lemma whose source label is \texttt{lem:orth} in the article (source0.tex lines 612--617, source label \texttt{lem:orth}), delivered together with the article's own complete original proof of it (source0.tex lines 619--663, one \texttt{proof} environment of 45 lines). No mathematics is written, repaired or re-derived: statement and proof are the article's own text line for line. Required context retained uncounted: prop:dBshift (lines 151--161); cor:split (lines 586--592). Every definition, display and label of those lines is retained unchanged. Omitted from this package and not redistributed: 1--150, 162--585, 593--611, 618--618, 664--960. Nothing inside a retained excerpt is omitted or altered: every retained body is delivered exactly as the article's lines are written, so no omission receipt of any kind accompanies this package. Citations: the retained excerpts carry no citation command at all, so no bibliography entry of the article is transcribed and no reference is redistributed. Reference adaptation: none was needed and none was made; every reference inside the delivered unit resolves inside it, so no unresolved reference or citation is produced. Printed slips kept literal: no printed word, symbol, hypothesis or number of the counted statement or of its proof was altered. Layout and class: the article's own class is \texttt{amsart} with packages including mathtools, amssymb, hyperref; the wrapper is the lane's pinned \texttt{amsart} editorial wrapper at 10pt on A4, which restates the theorem-like environments the retained text needs and adds the article's own macro definitions that the retained text needs (none), with extra wrapper packages (bm, bbm, dsfont, xspace, cancel, mathtools). The delivered numbering is the unit's own. Assets: no retained span contains a figure environment, a graphics-inclusion command, a PDF-page inclusion, a table or a TikZ picture; no image file is copied into this package, the package declares no asset dependency and no image file is redistributed with it. Licence: version arXiv:2508.08968v1 of this article is distributed under the Creative Commons Attribution 4.0 International licence, as recorded in the arXiv licence metadata for this exact version and shown on the abstract page https://arxiv.org/abs/2508.08968v1. A full rights notice is delivered at licenses/arxiv-2508.08968-CC-BY-4.0.txt. Characters: every retained excerpt is pure ASCII, so the delivered engine, XeLaTeX with its default Latin Modern text font, reports no missing character.
\begin{proposition}
\label{prop:dBshift}
For a subset $B$ of the label set $A$ and $\omega_{|B} \in \{0,1\}^B$, we have 
\renewcommand{\labelenumi}{\roman{enumi})}
\begin{enumerate}
    
\item $ C_{s_{B}^{\omega_{|B}}} \left(\mathcal{P}^{A,l}(G)\right) = \mathcal{P}^{A\setminus B,l}(G)$,

\item $ \delta_B \left(\mathcal{P}^{A,l}(G)\right) = \mathcal{P}^{A\setminus B,l -|B|}(G)$.
\end{enumerate}
\end{proposition}

\begin{corollary}
\label{cor:split}
For a character $A$-cube  $P : \{0,1\}^A \to \widehat{G}$, the value of the inner product $\left\langle P \right\rangle_{A,l}$ can be expressed as
$$ \left\langle P \right\rangle_{A,l}  = 
\left\langle \delta_i P \right\rangle_{A\setminus \{i\},l} 
\left\langle P \circ s_i^0 \right\rangle_{A \setminus \{i\},l-1}.$$
\end{corollary}

\begin{lemma}
\label{lem:orth}
 {\it (Orthogonality lemma)}
Let $P$ be an  $A$-cube of characters on $G$ of degree $k\in \{-1, 0, \dots, d \}$. Then $P$  satisfies the following orthonormality condition with respect to the $2^d$-ary inner product of degree $l\in \{-1, 0, \dots, d \}$:
$$ \langle P \rangle_{A,l} = \begin{cases} 1, & \text{if } k + l < d, \\ 0, & \text{if } k + l \ge d. \end{cases}$$
\end{lemma}

\begin{proof}
We will proceed by induction on the dimension  $d = |A|$. The case $A=\emptyset$ is left as an exercise, and we will discuss  the situation for 
 $d=1$ in more detail.  
 
Consider  the set of $1$-dimensional cubes in $\widehat{G}$, which are pairs of characters on $G$. This set is subject to a a filtration $\{I\}=\mathcal{P}^{A,-1}(\widehat{G}) \subset \mathcal{P}^{A,0}(\widehat{G}) \subset \mathcal{P}^{A,1}(\widehat{G}) = \widehat{G}^{\{0,1\}^{A}}$.

The set $\mathcal{P}^{A,-1}(\widehat{G})$ of cubes of degree $-1$ contains only one pair of trivial characters $(1_0,1_1)$. We can verify that for all possible choices of $l \in \{-1,0,1\}$, the claim holds:
$$ \langle \left(1_0,1_1\right) \rangle_{A,-1} = 1_0(0)\overline{1_1(0)} = 1,$$
$$ \langle \left(1_0,1_1\right) \rangle_{A,0} = \mathbb{E}_{x \in G}1_0(x)\overline{1_1(x)} = 1,$$
$$ \langle \left(1_0,1_1\right) \rangle_{A,1} = \mathbb{E}_{x_0,x_1 \in G}1_0(x_0)\overline{1_1(x_1)} = \mathbb{E}_{x \in G}1_0(x) \, \overline{\mathbb{E}_{x \in G}1_1(x)}= 1.$$

An $A$-cube $P$ is of degree 0  whenever it is a nontrivial constant map, i.e., $\chi_0 = \chi_1 \neq 1$. 
In this case, the claim follows from the fact that nontrivial characters have a unit $L_2$ norm and average to zero over $G$:
$$ \langle \left(\chi,\chi\right) \rangle_{A,-1} = \chi(0)\overline{\chi(0)} = 1,$$
$$ \langle \left(\chi, \chi\right) \rangle_{A,0} = \mathbb{E}_{x \in G}\chi(x)\overline{\chi(x)} = 1,$$
$$ \langle \left(\chi,\chi\right) \rangle_{A,1} = \mathbb{E}_{x_0,x_1 \in G}\chi(x_0)\overline{\chi(x_1)} = \mathbb{E}_{x \in G}\chi(x) \, \overline{\mathbb{E}_{x \in G}\chi(x)}= 0.$$

The case $k=1$ corresponds to pairs $P=(\chi_0, \chi_1)$ where $\chi_0 \neq \chi_1$. For such pairs, the claim follows from the orthogonality of characters in the standard Hermitian inner product on $\mathbb{C}^G$. Specifically, at least one character in a non-constant pair must average to zero over $G$:
$$\langle (\chi_0,\chi_1)\rangle_{A,-1} =  \chi_0(0)\overline{\chi_1(0) }= 1, $$
$$\langle (\chi_0,\chi_1)\rangle_{A,0} = \mathbb{E}_{x\in G} \chi_0(x)\overline{\chi_1(x) }= 0, $$
$$ \langle \left(\chi_0,\chi_1\right) \rangle_{A,1} = \mathbb{E}_{x_0,x_1 \in G}\chi_0(x_0)\overline{\chi_1(x_1)} = \mathbb{E}_{x \in G}\chi_0(x) \, \overline{\mathbb{E}_{x \in G}\chi_1(x)}= 0.$$

Similarly, for  $|A|=d>1$ and $l=-1$,  we  observe that $\langle P \rangle_{A,-1} =1$ for all $P \in \widehat{G}^{\{0,1\}^{A}}$. This follows because when $l = -1$, the multiplicative discrete derivative $\overline{\Delta}^A_{p,-1}(P)$ is evaluated solely at the $2^d$-dimensional zero point $p =0 \in \mathcal{P}^{A,-1}( G)$. 

On the other extreme, when $l=d$, we average over the entire space ${G}^{{\{0,1\}}^A}$ and can interchange the order of averaging and multiplication  
$$\langle P \rangle_{A,d} =
 \mathbb{E}_{p \in {G}^{{\{0,1\}}^A}}\prod_{\omega \in \{0,1\}^A} \mathcal{C}^{|\omega|} \chi_\omega (x_\omega) = 
\prod_{\omega \in \{0,1\}^A} \mathcal{C}^{|\omega|} \mathbb{E}_{x \in G}\chi_\omega (x).$$ 
This expression is nonzero if and only if $P$ is the $A$-cube of trivial characters  $I \in \mathcal{P}^{A,-1}(\widehat{G})$, i.e., has degree $k=-1$. In this case, the average equals $1$.

For $k=-1$, all multiplicative discrete derivatives of the trivial $A$-cube of characters $I$ are equal to one. Hence, after averaging over any $\mathcal{P}^{A,l}(G)$, we get  
$$\langle I \rangle_{A,l} = \mathbb{E}_{p \in  \mathcal{P}^{A,l}(G)} \overline{\Delta}^A_{p,l} (I) = 1.$$

Now, we may assume that $d\ge 2$, $k\in \{0,1, \dots d\} $,  $l\in \{0,1, \dots d-1\}$ and proceed with the induction step. 

Consider an $A$-cube of characters $P$ of degree $k$. Then, by Proposition \ref{prop:dBshift}, $\delta_i P$ is an $A\backslash  \{i\}$-cube of characters of degree $k-1$ and $ P \circ s_i^0$ is an $A\backslash \{i\}$-cube of characters of degree $k$. By induction hypothesis, we have
$$\left\langle \delta_i P \right\rangle_{A\setminus \{i\},l} = \begin{cases} 1, & \text{if } k-1 +l  < d -1, \\ 0, & \text{otherwise,} \end{cases},$$
and
$$\left\langle P \circ s_i^0 \right\rangle_{A\setminus \{i\},l-1} = \begin{cases} 1, & \text{if } k +l -1 < d -1, \\ 0, & \text{otherwise.} \end{cases}$$
Using Corollary \ref{cor:split}, it follows 
$$ \left\langle P \right\rangle_{A,l}  = 
\left\langle \delta_i P \right\rangle_{A\setminus \{i\},l} 
\left\langle P \circ s_i^0 \right\rangle_{A \setminus \{i\},l-1} =\begin{cases} 1, & \text{if } k + l  < d, \\ 0, & \text{otherwise.} \end{cases}
$$
\end{proof}

\end{document}
